\begin{abstract}
A recent proposal governs cloud data pipelines with four policy-gated LLM agents and reports
substantial gains over static orchestration. The original study does not give enough implementation
detail to pin down the exact model-in-the-loop evaluation condition behind those gains. We build the
system and test it with a real, remote language model, in a \NRunsTotal{}-run campaign that crosses
four injected fault scenarios with eight configurations (three non-agent comparators, the full
system, and four single-agent ablations) and splits into live-agent, non-agent, and mock-model arms.
Under a live model, the full system changes mean time to recovery (MTTR) by \NRelpctMttrSFull{}\% and
cost by \NRelpctCostUnitsFull{}\% relative to static orchestration, yet it \emph{increases} manual
interventions by \NRelpctManualInterventionsFull{}\%. Two idealized non-LLM baselines, each given a
fixed modelled remediation latency, recover faster still, changing MTTR by
\NRelpctMttrSRuleBased{}\% and \NRelpctMttrSAutoscale{}\%. Running the same \texttt{full}
configuration under a deterministic mock, instead of the live model, drops recovery time to
\NLiveVsMockMttrSMock{}\,s, about three orders of magnitude below the live figure, so a mock
evaluation and a live evaluation of one system point to different conclusions. Grading a generated
adversarial corpus of $\NAdversarialCases{}$ cases against an oracle independent of the policy code
gives containment $\geq\NAdversarialWilsonLo{}$, but only after fixing two contract-layer gaps the
corpus itself exposed. The evidence is bounded: a single host, synthetic data with injected faults, a
simulated on-call engineer, one live model family, and baselines that are modelled rather than
deployed.
\end{abstract}

\begin{IEEEkeywords}
LLM agents, cloud data pipeline governance, policy-as-code, Open Policy Agent, AIOps,
evaluation methodology, adversarial evaluation, mock-vs-live evaluation validity
\end{IEEEkeywords}
